% theory_reliability_foundations.tex —— 可靠性评估通用理论：元件模型、网络约简、充裕度与蒙特卡洛
% 本文件为理论层章节，遵循 docs/modules/THEORY_WRITING_GUIDE.md。

\section{电力系统可靠性评估的数学基础}\label{sec:rel-th}

\begin{theorynote}
本章为通用数学理论，按电力系统可靠性工程的标准体系（Billinton--Allan~\cite{rel:billinton-allan}）撰写，
覆盖：两状态元件的连续时间马尔可夫（CTMC）可用度模型、串并联网络的频率—持续时间（F\&D）约简与
最小割集、发电充裕度的停运容量概率表（COPT）递推与 LOLP/LOLE/LOLF/LOLD 指标、非序贯与序贯蒙特卡洛
估计量及其无偏性与变异系数（CoV）收敛率、IEEE Std 1366 配电可靠性指标（SAIFI/SAIDI/CAIDI/ASAI）、
以及尾部风险 VaR/CVaR。全部结论以平台实现（\srcpath{src/reliability/}）为落点，规模上限、保真度与
适用范围通过结果范围和有效性标志声明。纯交流系统采用交流网络 DC-OPF 削减口径
（\texttt{model\_scope="ac-only-dcopf"}）；含直流网络的系统采用
\texttt{hybrid-acdc-network-lp}，联合计入 AC/DC 有功平衡、直流负荷削减及 VSC/DC-DC 有界传输。
两种口径均不认证非线性交流无功与电压可行性。
\end{theorynote}

\subsection{记号与约定}\label{sec:rel-th-notation}
\begin{itemize}
  \item $\lambda$：元件故障率（占/年，occ/yr）；$\mu=1/r$：修复率；$r$：平均修复时间 MTTR（yr 或 hr）；
        $m=1/\lambda$：平均无故障时间 MTTF；$T=m+r$：平均循环时间（MTBF）；
  \item $A$：可用度（稳态在运行概率）；$U=1-A$：强迫停运率 / 不可用度（FOR）；$f$：状态频率（occ/yr）；
  \item $H=8760$（或 $8736$）：年小时数；$L$：负荷（MW）；$C$：可用发电容量（MW）；
  \item EENS：期望缺供电量（MWh/yr）；EDNS：期望缺供功率（MW）；
  \item LOLP/LOLE/LOLF/LOLD：失负荷概率/期望（hr/yr）/频率（occ/yr）/持续（hr/occ）；
  \item 主动型元件：$p_d$ 每次需求失效概率、$\nu_d$ 需求频率（次/年），等效年频率 $\lambda_a=\nu_d p_d$。
\end{itemize}
除特别声明，所有元件相互独立，故障/修复时间服从指数分布（马尔可夫假设）。

\subsection{两状态元件模型}\label{sec:rel-th-twostate}\label{sec:rel-param}
可修元件在“运行”（up）与“故障”（down）间交替，构成两状态连续时间马尔可夫链，生成元矩阵
\begin{equation}
 Q=\begin{bmatrix}-\lambda & \lambda\\ \mu & -\mu\end{bmatrix}.
 \label{eq:rel-th-Q}
\end{equation}

\begin{theorem}[稳态可用度与频率]\label{thm:rel-th-avail}
两状态元件~\eqref{eq:rel-th-Q} 的稳态概率、频率为
\begin{equation}
 A=\frac{\mu}{\lambda+\mu}=\frac{m}{m+r},\qquad
 U=\frac{\lambda}{\lambda+\mu}=\frac{\lambda r}{1+\lambda r},\qquad
 f=\lambda A=\mu U=\frac{1}{m+r}.
 \label{eq:rel-th-AUf}
\end{equation}
\end{theorem}

\begin{proof}
稳态分布 $\pi=[A,U]$ 满足 $\pi Q=0$ 与 $A+U=1$。$\pi Q=0$ 的首列给出平衡方程
$-\lambda A+\mu U=0$，即 $\lambda A=\mu U$；代入 $A+U=1$ 解得 $A=\mu/(\lambda+\mu)$、
$U=\lambda/(\lambda+\mu)$。以 $\mu=1/r$ 代入得 $U=\lambda/(\lambda+1/r)=\lambda r/(1+\lambda r)$。
状态频率为“离开该状态的稳态流率”$f=\lambda A=\mu U$；又 $\lambda A=\lambda\mu/(\lambda+\mu)
\allowbreak=1/(1/\lambda+1/\mu)\allowbreak=1/(m+r)$。
\end{proof}

平台以单一\emph{解析器}把异构输入（$\lambda$、MTTR、MTBF、FOR、$p_d/\nu_d$）折算为
式~\eqref{eq:rel-th-AUf} 的规范量 $(\lambda,r,U)$，供所有方法共享，并记录 MTBF 语义与数据来源
（\srcpath{include/hacdcpf/reliability/reliability\_assessment.hpp:resolve\_reliability\_params}）。

\subsection{网络约简：串并联与最小割集}\label{sec:rel-th-network}
独立两状态 CTMC 的串并联子网可用稳态可用度、边界穿越频率和平均故障持续时间逐级合并。下式同时列出
工程中常用的小不可用度（$U_i\ll1$）近似；平台实现保留精确稳态乘积与边界穿越频率，而非只计算一阶式。

\begin{proposition}[串联合并]\label{prop:rel-th-series}
$n$ 个串联元件（任一故障即系统故障），系统三元组为
\begin{equation}
 U_s=1-\prod_i(1-U_i)\approx\sum_iU_i,\qquad
 \lambda_s=\sum_i\lambda_i,\qquad
 r_s=\frac{\sum_i\lambda_ir_i}{\sum_i\lambda_i}.
 \label{eq:rel-th-series}
\end{equation}
\end{proposition}

\begin{proof}
系统可用要求全部元件可用，$A_s=\prod_iA_i$，故 $U_s=1-\prod_i(1-U_i)$，一阶展开得
$\sum_iU_i$。系统故障频率为各元件“单独触发系统故障”的频率之和
$f_s=\sum_i f_i\prod_{j\ne i}A_j\approx\sum_i\lambda_i$（$A_j\approx1$，$f_i\approx\lambda_i$）；
由 $\lambda_s=f_s/A_s\approx f_s$ 得 $\lambda_s=\sum_i\lambda_i$。等效修复时间由 $U_s=\lambda_s r_s$
（小量下 $U\approx\lambda r$）与 $U_s\approx\sum\lambda_ir_i$ 联立得 $r_s=\sum\lambda_ir_i/\sum\lambda_i$。
\end{proof}

\begin{proposition}[双元件并联合并]\label{prop:rel-th-parallel}
两个并联元件（同时故障才系统故障），在 $\lambda_ir_i\ll1$ 下
\begin{equation}
 U_p=U_1U_2\approx\lambda_1r_1\lambda_2r_2,\qquad
 \lambda_p\approx\lambda_1\lambda_2(r_1+r_2),\qquad
 r_p=\frac{r_1r_2}{r_1+r_2}.
 \label{eq:rel-th-parallel}
\end{equation}
\end{proposition}

\begin{proof}
独立性给 $U_p=U_1U_2$。系统进入故障态需一元件先故障（频率 $\approx\lambda_1$ 或 $\lambda_2$）而另一
元件恰处于故障态（概率 $U_2$ 或 $U_1$）：$f_p\approx\lambda_1U_2+\lambda_2U_1
=\lambda_1\lambda_2 r_2+\lambda_2\lambda_1 r_1=\lambda_1\lambda_2(r_1+r_2)$。等效
$\lambda_p=f_p/A_p\approx f_p$。$r_p$ 由 $U_p=\lambda_p r_p$ 得
$r_p=U_1U_2/[\lambda_1\lambda_2(r_1+r_2)]=r_1r_2/(r_1+r_2)$。
\end{proof}

一般（含网络约束、非串并联）物理拓扑用\emph{最小割集}表述：系统故障事件
$=\bigcup_k(\text{割集 }K_k\text{ 全故障})$，一阶近似 $U_{\mathrm{sys}}\approx\sum_k\prod_{i\in K_k}U_i$，
主导项为低阶（一阶、二阶）割集。平台对一般\emph{物理网络}仍以蒙特卡洛状态抽样统计削减（见
\S\ref{sec:rel-th-mc}）；对显式信息服务成功路径则已实现共享元件保持、QoS 筛选、最小割集和精确容斥
（\S\ref{sec:rel-icp-avail}），两者不可混称。

\subsection{发电充裕度：停运容量概率表与 LOLE/LOLF}\label{sec:rel-th-copt}
发电充裕度评估把机组停运卷积成\emph{停运容量概率表}（COPT）。

\begin{proposition}[COPT 递推]\label{prop:rel-th-copt}
设已叠加若干机组得累积停运概率 $P(X)=\Pr\{\text{停运容量}\ge X\}$。再叠加一台容量 $C$、强迫停运率
$U$ 的机组，则
\begin{equation}
 P'(X)=(1-U)\,P(X)+U\,P(X-C).
 \label{eq:rel-th-copt}
\end{equation}
频率表按同样卷积递推（附加 $\pm$ 机组故障/修复率的频率增量）。
\end{proposition}

\begin{proof}
新机组以概率 $1-U$ 在运（停运容量不变，贡献 $P(X)$），以概率 $U$ 停运（停运容量增加 $C$，
把原表右移 $C$，贡献 $P(X-C)$）。两互斥事件相加即~\eqref{eq:rel-th-copt}。
\end{proof}

由 COPT 与负荷 $L$，失负荷概率与期望为
\begin{equation}
 \mathrm{LOLP}=\Pr\{C<L\}=P(\,\text{停运}>C_{\mathrm{inst}}-L\,),\qquad
 \mathrm{LOLE}=\sum_{t}\mathrm{LOLP}(L_t)\,\Delta t,
 \label{eq:rel-th-lole}
\end{equation}
$C_{\mathrm{inst}}$ 为装机容量。F\&D 法进一步给出失负荷频率 $\mathrm{LOLF}$（越过 $L$ 的净向上频率）与
持续 $\mathrm{LOLD}=\mathrm{LOLE}/\mathrm{LOLF}$。平台以 F\&D 结构体
\srcpath{FrequencyDurationResult}（\fld{capacity\_outage\_levels}、\fld{cumulative\_probability}、
\fld{cumulative\_frequency}、\fld{lolp}、\fld{lole\_fd}、\fld{lolf\_fd}、\fld{lold}）承载 COPT 与指标。

\begin{example}[两机组 COPT 与充裕度指标]\label{ex:rel-th-copt}
两台 $60$~MW 机组，强迫停运率 $q=0.05$（$p=0.95$），$\mathrm{MTTF}=1900$~h、$r=100$~h
（校验 $U=r/(\mathrm{MTTF}{+}r)=0.05$，$\lambda=8760/1900=4.61/\text{yr}$、$\mu=8760/100=87.6/\text{yr}$）。
停运容量 $X$ 的概率表为 $P(0){=}p^2{=}0.9025$、$P(60){=}2pq{=}0.095$、$P(120){=}q^2{=}0.0025$。
设峰荷 $L=50$~MW、总容量 $120$~MW，则失负荷当可用容量 $120-X<50$，即 $X\ge120$（两机全停）：
\begin{equation*}
 \mathrm{LOLP}=P(X{\ge}120)=0.0025,\qquad \mathrm{LOLE}=8760\times0.0025=21.9~\text{h/yr}.
\end{equation*}
两机全停态的离开率为 $2\mu=175.2/\text{yr}$，故 $\mathrm{LOLF}=0.0025\times175.2=0.438~\text{occ/yr}$、
$\mathrm{LOLD}=\mathrm{LOLE}/\mathrm{LOLF}=50$~h（$=1/(2\mu)$，两机中任一先修复即结束）。若峰荷升至 $70$~MW，
则 $X{\ge}60$ 即失负荷，$\mathrm{LOLP}=0.095{+}0.0025=0.0975$——凸显充裕度对备用裕度的强敏感性。
\end{example}

\subsection{公共后果模型：突变算子、最小切负荷与引擎阶梯}\label{sec:rel-th-conseq}
失效模式\emph{不}直接等于元件停运，而经\emph{后果算子}映射为网络 / 控制\emph{突变}：
\begin{equation}
 G_m=\Phi_m(G_0,x_m),
\end{equation}
$G_0$ 为健康富系统、$x_m$ 为模式严重度、$\Phi_m$ 产生一或多个突变：拓扑（强迫停运 / 卡开 / 强迫切负荷）、
容量（按残余容量因子降额）、控制（定值冻结 / 失控 / 失去构网）、保护（拒动后区域扩展）、观测（量测 /
通信退化）、恢复（禁止某动作，如拒合）。硬停运在同一目标上压制降额与失控；不兼容的强迫开 / 合突变报为
冲突；所选物理引擎无法表示的突变标为\emph{不支持}而非人为置零。对失效态 $x$，各方法最终求\emph{最小切负荷}
\begin{equation}
 S(x)=\min_{y\in\mathcal F(G_x)}\sum_iw_ip_i^{sh},
 \label{eq:rel-th-minshed}
\end{equation}
$\mathcal F(G_x)$ 为所选物理模型的可行运行 / 恢复域。实现把 VOLL 置于可行发电边际成本之上，使源成本仅在
最小切负荷解间打破简并；可靠性只读 $S$ 与 $s_i$，不读经济目标值。三级\emph{物理后果引擎阶梯}为：
\begin{center}\footnotesize
\begin{tabular}{@{}L{2.9cm}L{3.4cm}L{7.3cm}@{}}
\toprule
引擎 & 使用者 & 主物理 / 有效边界\\
\midrule
交流 DC-OPF & 纯交流 MC/FMEA & $P{=}B\theta$、节点平衡、源 / 支路限、切负荷；无电压幅值 / 无功可行性\\
混合 AC/DC 网络 LP & 混合 MC、元件 / 故障模式 FMEA & AC 角—流、DC 平衡、有界 VSC/DC-DC 传输、储能 / DER；DC 传输线性稳态，无非线性 AC 校验\\
耦合 LinDistFlow 恢复 MILP & 三阶段模型 & AC 有 / 无功 $+$ DC 有功平衡、平方电压、辐射森林、开关、双向换流、DER / 储能时序；线性化损耗，拒 LCC 与多端口路由器\\
\bottomrule
\end{tabular}
\end{center}
依据：\srcpath{src/reliability/reliability\_assessment.cpp:evaluate\_failed\_network\_state}。

\subsection{蒙特卡洛可靠性评估}\label{sec:rel-th-mc}
当网络约束、多重故障与调度重分布使解析枚举不可行时，蒙特卡洛（MC）以状态抽样估计指标。

\paragraph{非序贯（状态抽样）。}
每次抽样对各元件独立取 Bernoulli 状态 $x_i\sim\mathrm{Bern}(U_i)$（$1$=故障），对该系统状态解
DC-OPF 得最小削减 $C(\mathbf x)$（MW），估计
\begin{equation}
 \widehat{\mathrm{EENS}}=\frac1N\sum_{k=1}^N C(\mathbf x^{(k)})\cdot H,\qquad
 \widehat{\mathrm{EDNS}}=\frac1N\sum_{k=1}^N C(\mathbf x^{(k)}).
 \label{eq:rel-th-nsq}
\end{equation}

\begin{theorem}[无偏性与 CoV 收敛]\label{thm:rel-th-cov}
估计量~\eqref{eq:rel-th-nsq} 无偏：$\mathbb E[\widehat{\mathrm{EDNS}}]=\mathbb E[C(\mathbf x)]=\mathrm{EDNS}$；
样本均值的变异系数
\begin{equation}
 \mathrm{CoV}=\frac{\sigma_{\widehat{}}}{\widehat\mu}=\frac{\sigma_C}{\mu_C\sqrt N}=O\!\left(N^{-1/2}\right),
 \label{eq:rel-th-cov}
\end{equation}
$\sigma_C,\mu_C$ 为单样本削减的标准差与均值。停机准则为 $\mathrm{CoV}\le\varepsilon$。
\end{theorem}

\begin{proof}
$\mathbf x^{(k)}$ 独立同分布，$\mathbb E[\tfrac1N\sum_kC(\mathbf x^{(k)})]=\mathbb E[C(\mathbf x)]$，
无偏。样本均值方差 $\mathrm{Var}(\widehat\mu)=\sigma_C^2/N$，标准差 $\sigma_C/\sqrt N$；
除以均值 $\mu_C$ 即 $\mathrm{CoV}=\sigma_C/(\mu_C\sqrt N)$，随 $N$ 以 $N^{-1/2}$ 衰减。
\end{proof}

CoV 的 $O(N^{-1/2})$ 率与\emph{维数无关}，是 MC 相对枚举的根本优势；但稀有事件（$\mu_C$ 小、
$\sigma_C/\mu_C$ 大）收敛慢，需重要抽样等方差缩减。依据：\srcpath{src/reliability/reliability\_assessment.cpp:run\_nonsequential\_mc}
（\fld{cov\_threshold}、\fld{max\_iterations}、\fld{final\_cov}）。

\paragraph{序贯（时序模拟）。}
按时间推进：各元件抽指数分布的运行/修复时长（率 $\lambda_i$/$\mu_i$）生成时序状态，逐时对
含负荷曲线的系统解 DC-OPF，累计年度 EENS 与\emph{越限次数}得 LOLF。序贯法额外给出频率与持续
（LOLF、LOLD）及其年际分布，代价是每样本为完整一年。依据：
\srcpath{src/reliability/reliability\_assessment.cpp:run\_sequential\_mc}、
\srcpath{build\_ieee\_rts24\_load\_profile}。

\subsection{配电可靠性指标（IEEE Std 1366）}\label{sec:rel-th-1366}
以负荷点 $i$ 的故障率 $\lambda_i$、年停电时间 $U_i$（hr/yr）与客户数 $N_i$ 定义系统指标：
\begin{align}
 \mathrm{SAIFI}&=\frac{\sum_i\lambda_iN_i}{\sum_iN_i},&
 \mathrm{SAIDI}&=\frac{\sum_iU_iN_i}{\sum_iN_i},&
 \mathrm{CAIDI}&=\frac{\mathrm{SAIDI}}{\mathrm{SAIFI}},\\
 \mathrm{ASAI}&=1-\frac{\mathrm{SAIDI}}{H},&
 \mathrm{ASUI}&=1-\mathrm{ASAI}.&&
\end{align}
SAIFI 为户均停电次数、SAIDI 为户均停电时长、CAIDI 为次均停电时长、ASAI 为供电可用率。平台以
\srcpath{DistributionIndices}（\fld{saifi}/\fld{saidi}/\fld{caidi}/\fld{asai}/\fld{asui} 及逐节点
\fld{nodal\_cif}/\fld{nodal\_cid}）承载，需负荷点客户数据方计算（\fld{compute\_distribution\_indices}）。

\subsection{尾部风险：VaR 与 CVaR}\label{sec:rel-th-var}
可靠性关注\emph{尾部}的高损失场景。对年损失 $E$（如年 EENS）与置信 $\alpha$，
$\mathrm{VaR}_\alpha(E)=\inf\{e:\Pr[E\le e]\ge\alpha\}$，严格的期望损失（Expected Shortfall）为
\begin{equation}
 \mathrm{ES}_\alpha(E)=\frac1{1-\alpha}\int_\alpha^1\mathrm{VaR}_u(E)\,\mathrm du
 =\min_{\eta\in\mathbb R}\Big\{\eta+\frac1{1-\alpha}\mathbb E[(E-\eta)^+]\Big\},
 \label{eq:rel-th-cvar}
\end{equation}
右式为 Rockafellar--Uryasev 表征~\cite{rel:rockafellar}；连续分布下等于
$\mathbb E[E\mid E\ge\mathrm{VaR}_\alpha]$。CVaR 是\emph{相干}风险度量（次可加、单调、平移不变、正齐次），
优于 VaR。\textbf{实现注}：平台按经验分布对最坏 $1-\alpha$ 概率质量积分；当 VaR 处有并列样本时只取补足
尾概率所需的分数质量，因此与上式的经验 ES 一致，而不是把全部 VaR 原子重复计入。序贯年样本为首选输入。
承载于 \srcpath{TailRiskMetrics}
（\fld{eens\_var}/\fld{eens\_cvar}/\fld{eens\_percentiles}）。

\subsection{采样不确定性与稀有事件}\label{sec:rel-th-uncert}
对 $n$ 个独立状态样本，缺供功率估计的标准误与渐近 $1-\alpha$ 区间为
\begin{equation}
 \mathrm{SE}(\widehat{\mathrm{EDNS}})=\frac{\widehat\sigma_S}{\sqrt n},\qquad
 \bar S\pm z_{1-\alpha/2}\frac{\widehat\sigma_S}{\sqrt n},\qquad
 \mathrm{SE}(\widehat{\mathrm{PLC}})=\sqrt{\frac{\widehat p(1-\widehat p)}{n}}.
\end{equation}
稀有失负荷宜用 Wilson 或精确二项区间；序贯年样本须以\emph{年级}方差 / 区间而非把小时当独立。极小
PLC/LOLE 下朴素 MC 统计低效。平台已实现优势比扭曲重要抽样、似然比无偏修正和 ESS 诊断；交叉熵、子集模拟与
按故障阶分层不属于该接口范围，完整推导见 \S\ref{sec:rel-exact}。这与 CoV 停机准则（定理~\ref{thm:rel-th-cov}）
互补：CoV 度量\emph{相对}收敛，SE / 区间度量\emph{绝对}不确定性。

\subsection{薄弱环节辨识（后处理）}\label{sec:rel-th-weakpoint}
评估后按元件对\emph{失负荷态}的贡献排序，定位薄弱环节。平台的元件重要度定义为
\srcpath{ReliabilityResult::ComponentImportance}，是\emph{共现 / 归因}度量，而非 Birnbaum 边际重要度
$\partial\mathrm{EENS}/\partial U_c$：以“元件 $c$ 停运时观测到的失负荷削减份额”定义
\begin{align}
 \text{loss\_weighted\_risk}_c&=\frac{\sum_k\lambda_k\,\mathbf 1[c\in\text{down}(k)]\,S_k}{\sum_k\lambda_k S_k},\\
 \Pr(c\text{ 停}\mid\text{系统失负荷})&=\frac{\sum_k\lambda_k\,\mathbf 1[c\in\text{down}(k)]\,\mathbf 1[S_k>\varepsilon]}{\sum_k\lambda_k\,\mathbf 1[S_k>\varepsilon]},
\end{align}
并给出共停关联 EENS $\text{associated\_eens}_c$。\textbf{诚实口径}：多重失效态把整段削减\emph{归因于每个}
停运元件，故各元件贡献之和可\emph{超过} 100\%——应读作\emph{排序}而非边际风险。精确状态接口另行计算
$\partial\mathrm{EENS}/\partial U_c$、Birnbaum 与 Fussell--Vesely，见 \S\ref{sec:rel-exact}。依据：\srcpath{src/reliability/reliability\_assessment.cpp}
（\fld{critical\_components}）。

\subsection{数值算例与基准（IEEE RTS-24）}\label{sec:rel-th-numexp}\label{sec:rel-th-state-consequence}
本节先规定 IEEE RTS-24 网络约束可靠性评估应求解的数学对象。该模型是后续代码修订与基准验收的
判据；在全部一致性条件通过前，不以某次蒙特卡洛汇总值反推模型正确。

\subsubsection{物理状态与唯一表示}
令 $\mathcal N$、$\mathcal E$、$\mathcal G$ 分别为 AC 母线、物理支路和机组集合，
$a_e(x),a_g(x)\in\{0,1\}$ 为状态 $x$ 下的可用标志。求解器使用的带电边集必须为
\begin{equation}
 \mathcal E(x)=\{e\in\mathcal E:a_e(x)=1\}.
 \label{eq:rel-rts-edge-set}
\end{equation}
若导入器为物理支路 $e$ 生成参数、控制或结果归因元数据 $t$，并以映射
$\rho(t)=e$ 关联，则 $t$ 不得再产生独立图边或独立随机变量：
\begin{equation}
 a_t(x)\equiv a_{\rho(t)}(x),\qquad
 \mathcal E_{\rm topo}(x)=\mathcal E_{\rm flow}(x)=\mathcal E(x).
 \label{eq:rel-rts-representation-invariant}
\end{equation}
式~\eqref{eq:rel-rts-representation-invariant} 是\emph{表示不变性}，不是数值近似。违反它会使拓扑筛选
与功率平衡看到两个不同系统，任何 EENS 结果均无物理解释。

由 $\mathcal E(x)$ 得连通分量 $\mathcal C(x)=\{C_1,\ldots,C_K\}$。按 MATPOWER DC 模型~\cite{rel:zimmerman}，对每个分量任取一个母线
$r_k\in C_k$ 固定 $\theta_{r_k}=0$，即可消除 DC 节点电纳矩阵的一维零空间：
\begin{equation}
 \operatorname{rank} B_k=|C_k|-1,\qquad \theta_{r_k}=0.
 \label{eq:rel-rts-angle-reference}
\end{equation}
因此 \texttt{SLACK} 最多提供一个可复用的角度规范，不是供电能力；\texttt{NoSlack} 也不等于失电岛。
对式~\eqref{eq:rel-rts-dcopf}，无需判定输入数据中“有没有平衡节点”：每个 $C_k$ 任取一个
$r_k$ 即可。一个没有导入 slack 标签、但含在线可调机组的分量仍可求解；一个保留 slack 标签、但没有
可用注入能力的分量仍会缺供。这里必须区分两类问题：普通 DC 潮流需要指定承担功率不平衡的 slack
注入，而 DC-OPF 已在所有母线同时执行式~\eqref{eq:rel-rts-balance}，功率不平衡由优化变量分配，
角度参考不承担额外注入。

\subsubsection{同一状态上的最小切负荷 DC 模型}
对时段 $h$ 的毛需求 $d_{ih}\ge0$，定义负荷削减 $s_{ih}\in[0,d_{ih}]$、机组出力 $p_{gh}$、
支路潮流 $f_{eh}$ 和相角 $\theta_{ih}$。静态电源、可再生电源、光伏、放电储能以及负的作者负荷
统一聚合为状态相关固定正注入 $\bar p^{\rm fix}_{ih}(x)\ge0$，并定义固定电源削减
$p^{gc}_{ih}$。网络状态后果模型为
\begin{subequations}\label{eq:rel-rts-dcopf}
\begin{align}
 S(x,h)=\min\quad & \sum_{i\in\mathcal N}s_{ih}, \label{eq:rel-rts-dcopf-obj}\\
 \text{s.t.}\quad
&\sum_{g\in\mathcal G_i}p_{gh}+\bar p^{\rm fix}_{ih}(x)-p^{gc}_{ih}
  +s_{ih}-d_{ih}
   =\sum_{e\in\delta(i)}K_{ie}f_{eh}, &&i\in\mathcal N,\label{eq:rel-rts-balance}\\
&f_{eh}=a_e(x)b_e(\theta_{o(e)h}-\theta_{d(e)h}), &&e\in\mathcal E,\label{eq:rel-rts-flow}\\
&-a_e(x)\bar f_e\le f_{eh}\le a_e(x)\bar f_e, &&e\in\mathcal E,\label{eq:rel-rts-limit}\\
&0\le s_{ih}\le d_{ih},\quad
  0\le p^{gc}_{ih}\le\bar p^{\rm fix}_{ih}(x), &&i\in\mathcal N,\label{eq:rel-rts-shed}\\
 &\theta_{r_kh}=0, &&C_k\in\mathcal C(x).\label{eq:rel-rts-ref}
\end{align}
\end{subequations}
其中 $K$ 为支路—节点关联矩阵，$b_e$ 为 DC 电纳。当前可靠性后果模型采用等权 MW 削减，且与
平台 DC-OPF 一致，不计支路移相角。若分量 $C_k$ 没有任何可用供给，对式~\eqref{eq:rel-rts-balance} 在 $C_k$ 上求和，边界潮流
相消，立即得到
\begin{equation}
 \sum_{i\in C_k}s_{ih}=\sum_{i\in C_k}d_{ih}.
 \label{eq:rel-rts-dead-island}
\end{equation}
所以无源岛的“全负荷切除”应由同一个优化模型自然产生，不需要以 slack 标签或设备类型作启发式回退。
更强地说，在采用 HL-II 再调度语义~\eqref{eq:rel-rts-hlii-redispatch} 时，
\begin{equation}
 p=0,\qquad f=0,\qquad \theta=0,\qquad s=d,\qquad
 p^{gc}=\bar p^{\rm fix}
 \label{eq:rel-rts-constructive-point}
\end{equation}
总是一个可行点；即使整个系统没有在线机组，或某孤岛固定注入大于负荷，模型也应返回可行最优解，而不是以
$n_g=0$、供给过剩或 \texttt{NoSlack} 返回求解失败。固定注入与毛负荷必须分开聚合；若先合成净负荷再令
$s_i\le\max(d_i-\bar p_i^{\rm fix},0)$，式~\eqref{eq:rel-rts-constructive-point} 将被错误破坏。

目标~\eqref{eq:rel-rts-dcopf-obj} 按字典序实现。令 $S^\star$ 为第一级最优值，第二级求解
\begin{equation}
 \min\ \sum_g C_g(p_g)+\sum_i\gamma_i p^{gc}_i
 \quad\text{s.t. 式~\eqref{eq:rel-rts-dcopf} 且 }\sum_i s_i=S^\star,
 \label{eq:rel-rts-lexicographic-second}
\end{equation}
其中 $\gamma_i\ge0$ 为固定电源削减代价。单一有限 VOLL 加权只有在给出目标界和求解容差时才能证明等价；否则
“VOLL 很大”只是工程近似，必须在结果有效性中声明。

\subsubsection{机组下界的三种运行语义}
机组约束不能在未声明评估层级时统一改写。至少有以下三个互不等价的模型：

\paragraph{A. HL-I 发电充裕度。}
忽略网络，仅比较可用容量与负荷：
\begin{equation}
 0\le p_{gh}\le a_g(x)\bar P_g.
 \label{eq:rel-rts-hli-dispatch}
\end{equation}
这与 COPT 的“在线容量可从 0 使用到 $\bar P_g$”一致，不消费 MATPOWER 的经济调度 $P_g^{\min}$。

\paragraph{B. HL-II 网络约束充裕度。}
保留式~\eqref{eq:rel-rts-dcopf} 的网络约束，并仍采用
\begin{equation}
 0\le p_{gh}\le a_g(x)\bar P_g.
 \label{eq:rel-rts-hlii-redispatch}
\end{equation}
该口径假设故障后允许充裕度再调度，但不模拟开停机、最小开停时间与爬坡。它适合把网络约束增量与
经典 RTS COPT 比较，但不能冒充生产模拟。

\paragraph{C. 固定/时序开机组合的运行可靠性。}
若开机状态 $y_{gh}\in\{0,1\}$ 来自 UC 或生产模拟，则
\begin{equation}
 a_g(x)y_{gh}P_g^{\min}\le p_{gh}\le a_g(x)y_{gh}\bar P_g.
 \label{eq:rel-rts-commitment}
\end{equation}
此时低负荷岛可能满足 $\sum_gP_g^{\min}>\sum_i d_i$。模型必须明确选择以下机制之一：消费 UC 的
紧急停机决策、加入有物理含义的弃电/外送变量、或把该状态报告为\emph{运行约束不可行}；不得把这种
供给过剩转换成负荷不足。没有 $y_{gh}$、爬坡、最小开停时间和弃电语义时，直接保留正 $P_g^{\min}$
或直接把全部 $P_g^{\min}$ 清零都不能称为通用正确模型。

\begin{center}\footnotesize
\begin{tabular}{@{}L{3.2cm}L{3.7cm}L{3.7cm}L{3.7cm}@{}}
\toprule
口径 & 机组下界 & 可与经典 RTS 比较 & 必须报告的限制\\
\midrule
HL-I COPT & $0$ & 是，发电充裕度 & 无网络\\
HL-II 再调度 & $0$ & 可比较网络约束增量 & 无 UC/爬坡/电压无功\\
固定/时序组合 & $a_gyP^{\min}$ & 仅与同一 UC 轨迹比较 & 需弃电或紧急停机语义\\
\bottomrule
\end{tabular}
\end{center}

\subsubsection{物理后果、基线缺额与数值失败必须分离}
对求解成功状态，定义原始物理削减 $S^{\rm raw}(x,h)$。同负荷下健康状态为 $x_0$，则
\begin{equation}
 S^0(h)=S^{\rm raw}(x_0,h),\qquad
 S^{\rm inc}(x,h)=\max\{0,S^{\rm raw}(x,h)-S^0(h)\}.
 \label{eq:rel-rts-raw-incremental}
\end{equation}
标准 EENS 使用 $S^{\rm raw}$；$S^{\rm inc}$ 是故障增量诊断。若 $S^0(h)>0$，应判定基准系统、负荷
曲线或运行语义未闭环，并同时报告原始量与增量量，不能静默用增量量替换标准指标。

令 $q(x,h)=1$ 表示优化器取得满足可行性容差的解。数值失败 $q=0$ 不是物理状态，故主估计量只能对
已认证状态定义：
\begin{equation}
 \widehat{\mathrm{EENS}}_{\rm cert}
 =\frac{H}{N}\sum_{k:q_k=1}S^{\rm raw}(x^{(k)},h_k),\qquad
 p_{\rm fail}=\frac1N\sum_k\mathbf1[q_k=0].
 \label{eq:rel-rts-certified-eens}
\end{equation}
若为失败状态临时取 $S^{\rm UB}=\sum_i d_i$，它只能形成保守上界
\begin{equation}
 \widehat{\mathrm{EENS}}^{\rm UB}
 =\widehat{\mathrm{EENS}}_{\rm cert}
 +\frac{H}{N}\sum_{k:q_k=0}S^{\rm UB}_k,
 \label{eq:rel-rts-eens-upper-bound}
\end{equation}
不得把式~\eqref{eq:rel-rts-eens-upper-bound} 标成普通 EENS。每个失败状态至少记录：状态键、连通分量、
在线源与上下界、求解器状态、原始/对偶残差及后端链。只有证明失败源于模型物理不可行，且模型已包含
式~\eqref{eq:rel-rts-shed} 所需的松弛后，才能把它解释为系统后果。

\begin{theorynote}
本节契约已在 AC-only HL-II 与混合 AC/DC 有功网络两条代码路径闭环实现。AC-only 入口强制开启
毛负荷削减、固定正注入削减、$P_g^{\min}=0$ 和
两阶段字典序目标：第一阶段最小化 $\sum_i s_i$，第二阶段加入 $\sum_i s_i=S^\star$ 等式后
最小化 PWL 发电成本与 $p^{gc}$ 代价。混合 LP 同样执行两阶段目标，VSC/DC--DC 紧急再调度区间包含
零传输。两条路径均使用唯一物理边集和每个 AC 连通分量一个角度参考。DC-OPF 在规范投影坐标中对
等式、不等式、盒界和有限性作后验认证；若后端失败或证书超差，评估立即抛出含残差与后端链的错误，
该状态不写入 EENS。序贯 MC 的 \fld{annual\_eens} 与 \fld{eens\_mwh\_yr} 逐小时累加
$S^{\rm raw}$，\fld{baseline\_eens\_mwh\_yr} 单列 $S^0$，
\fld{incremental\_eens\_mwh\_yr} 单列 $S^{\rm inc}$；LOLE/LOLF 也按原始失负荷标志计算。
因此健康状态已有缺额时，不会再被静默扣除出标准指标。
\end{theorynote}

\subsubsection{诊断数据与方法对比}
在 IEEE 可靠性测试系统 RTS-24（$24$ 母线、$33$ 机组、$38$ 支路；\srcpath{data/case24\_ieee\_rts.m}
$+$ \srcpath{apply\_ieee24\_reliability\_data}）上运行两类 MC（\srcpath{tools/pf\_doc\_benchmark.cpp}
模式 \texttt{rel}，\texttt{ac-only-dcopf} 口径）：
可靠性数据按 \texttt{case24\_ieee\_rts.m} 的 33 行机组和 38 行支路逐行映射，并以
母线号、$P_g^{\max}$ 和支路端点签名校验。第 15 行 $P_g^{\max}=0$ 是同步调相机，仍保留其
记录身份但不提供有功容量。旧实现按母线给同一母线全部机组复制一个 FOR/MTTR，并把并联支路按
端点合并匹配；这会改变 COPT 和网络状态概率，是结果偏高/偏低均可能发生的数据构模错误。
\begin{center}
\resizebox{\textwidth}{!}{%
\begin{tabular}{@{}lccccc@{}}
\toprule
层级/方法 & 负荷口径 & CoV/年数 & EENS/(MWh$\cdot$yr$^{-1}$) & LOLE/(hr$\cdot$yr$^{-1}$) & 诊断\\
\midrule
HL-I 当前机组映射（COPT） & 峰值 $2850$~MW & 精确卷积 & 仅容量缺额 & 见 FD & 不含网络故障\\
非序贯 MC（逐台数据，诊断） & 固定峰值、8760 h & $0.0314/2\times10^4$ 态 & $123\,156.8$ & $738.5$ & 未达到停止条件\\
序贯 MC（逐台数据，诊断） & 8736 h 曲线 & $0.1152/400$ 年 & $1\,222.6$ & $10.54$ & 未达到停止条件\\
确定性 HL-II 扫描 & N-0/N-1/N-2 & $2557/2557$ 态 & -- & -- & 失败 0、残差 0\\
\bottomrule
\end{tabular}
}
\end{center}
两类 MC 数值不作为模型可行性证明：非序贯使用固定峰值，序贯使用时序曲线，且本次运行的 CoV
尚未达到停止阈值。当前模型先用确定性扫描证明每个 N-0/N-1/N-2 状态均可认证，再把 MC
作为概率估计。修正后每小时采用
\begin{equation}
 S^{\mathrm{inc}}_{h}=\max\{0,S(x_h,L_h)-S(x_0,L_h)\},\qquad
 \mathrm{EENS}^{\mathrm{inc}}=\frac1{N_y}\sum_y\sum_h S^{\mathrm{inc}}_{y,h}.
\end{equation}
其中 $x_0$ 是同一负荷倍率下的 N-0 状态。该量按式~\eqref{eq:rel-rts-raw-incremental} 只能称为
故障增量 EENS；标准原始 EENS 仍须另报。上表两次 MC 均为诊断运行，不能当作最终收敛基准；
最终基准必须在预设 CoV 停止条件满足且失败计数为零后发布。
这里的 MC \texttt{converged} 仅表示统计误差停止条件；它不能覆盖状态级求解证书。后者若为
\texttt{converged=false} 或后验可行性失败，属于构模/数值求解缺陷，绝不作为“尚未抽够样本”处理，
也不允许继续产生可靠性指标。
\begin{boundarynote}
经典 RTS-24 的 EENS $\sim1.2\times10^3$~MWh/yr、LOLE $\sim9.4$~hr/yr 通常是 HL-I 发电充裕度口径，仅对机组 COPT 卷积，不抽样输电支路。发电与网络风险一般重叠，不能直接作相互独立的加法分解。应在相同状态与负荷上计算反事实增量
\begin{equation}
 \Delta\mathrm{EENS}_{\rm network}
 =H\,\mathbb E\!\left[S_{\rm HL\text{-}II}(x,h)-S_{\rm HL\text{-}I}(x,h)\right],
 \label{eq:rel-rts-network-increment}
\end{equation}
两模型采用相同机组可用状态和运行语义时，HL-II 可行域是 HL-I 可行域加网络限制后的子集，故该增量应非负。
无源岛削减与其余连通区域的 OPF 削减可按\emph{互斥母线集合}相加，但它们已经包含在
$S_{\rm HL\text{-}II}$ 内，不能再次叠加。式~\eqref{eq:rel-rts-eens-upper-bound} 只保留为历史
诊断定义；当前执行契约不接受任何数值失败状态，不能把失败状态的上界写入 EENS，也不能混入
物理网络风险。结果结构现已报告
\texttt{opf\_failed\_probability}、\texttt{dead\_island\_eens\_mwh\_yr} 和
\texttt{opf\_failed\_eens\_mwh\_yr}。当前闭环模型要求两个 \texttt{opf\_failed} 字段严格为零；
\texttt{dead\_island\_eens} 是求解成功后的物理归因，可以非零。历史失败回退值只保留为缺陷定位记录，
不能称为 RTS 物理风险；若状态求解返回 \texttt{converged=false} 或后验证书失败，入口立即抛出包含
状态键、求解器状态和残差的错误，调用方必须修复构模或数值算法后重新运行。

产生回退的具体原因是 MATPOWER 变压器同时存在两种表示：带变比支路保留在 \texttt{ac.branches}，并额外生成带
\texttt{source\_branch\_idx} 的 \texttt{Transformer2W} 元数据。旧拓扑图把两者都当成连通边，而 DC-OPF 只把 \texttt{ac.branches}
装入支路约束。因此旧实现中“变压器支路停运、元数据边仍在”的状态会被图判为有源连通岛，随后
OPF 面对实际断开的 $B$ 矩阵而失败；旧失败分支又把存活负荷全部计为削减。现行拓扑、AC-only DC-OPF、
混合 AC/DC LP 和随机元件集合均跳过链接元数据边，且失败状态直接抛错，不再生成 EENS 数值。

该故障链可写成同一状态向量上的两个不同边集：
\begin{equation}
 E_{\mathrm{graph}}=E_{\mathrm{branch}}\cup E_{\mathrm{transformer\ metadata}},\qquad
 E_{\mathrm{opf}}=E_{\mathrm{branch}} .
\end{equation}
若对应原始支路状态为 $x_b=1$ 而元数据仍为在线，旧实现得到
$G(E_{\mathrm{graph}})$ 连通、但 $B(E_{\mathrm{opf}})$ 已断开；拓扑预筛因此不切除该负荷，
而 OPF 无法为断开的实际分量建立可行功率平衡，触发“存活负荷全部削减”的回退。反向状态
（元数据停运而原始支路在线）则会采到一个不对应任何物理故障的随机事件。两种情况分别造成
假阴性和假阳性，不能用“保守模型”解释。现行图、OPF 和 MC 元件集合均跳过
\texttt{source\_branch\_idx>0} 的链接元数据，因此该双重表示不再产生独立随机状态。

表示层修订采用的不变量是：对任意链接元数据行 $t$，只要 $t.\texttt{source\_branch\_idx}>0$，它只能参与
结果归因和映射，不能参与图边、孤岛边或独立蒙特卡洛元件集合；其在线状态完全由对应
\texttt{ACBranch} 决定。该不变量由
\srcpath{tests/test\_reliability\_resolver.cpp} 中的两母线回归测试执行验证：健康状态拓扑连通，
停运原始支路后恰为两个 AC 岛、负荷岛状态为 \texttt{NoSlack}，且 DC-OPF 返回 10 MW 的可审计切负荷，
不再借助元数据边维持连通。自适应孤岛子系统也采用同一过滤规则（\srcpath{src/power\_flow/island\_detector.cpp}）。

同类但独立的错误来自把母线 \texttt{SLACK} 类型误作在线发电能力。进一步按
式~\eqref{eq:rel-rts-balance}--\eqref{eq:rel-rts-ref} 推导可知，修复并不是把判据改成“岛内是否有
在线源”，而是完全取消 DC-OPF 的平衡节点/在线源可行性预筛：无 slack 有源岛由局部角度参考求解，
有 slack 无源岛和无 slack 无源岛都由 $s=d$ 内生得到全切负荷。现行实现已经删除该在线源预筛，
并以无机组、固定注入过剩、断岛和链接变压器四类回归用例约束这一行为。
\end{boundarynote}

\subsection{与实现的对应}\label{sec:rel-th-impl}
\begin{center}\footnotesize
\begin{longtable}{@{}L{5.0cm}L{9.4cm}@{}}
\toprule
\textbf{理论对象} & \textbf{平台实现状态}\\
\midrule
\endfirsthead
\toprule
\textbf{理论对象} & \textbf{平台实现状态}\\
\midrule
\endhead
\bottomrule
\endfoot
两状态可用度~\eqref{eq:rel-th-AUf}、参数解析 &
  \implfull{src/reliability/reliability\_assessment.cpp:resolve\_reliability\_params}\\
非序贯（状态抽样）MC~\eqref{eq:rel-th-nsq}、CoV 停机 &
  \implfull{src/reliability/reliability\_assessment.cpp:run\_nonsequential\_mc}\\
序贯（时序）MC、LOLF/年际分布 &
  \implfull{src/reliability/reliability\_assessment.cpp:run\_sequential\_mc}\\
COPT 递推~\eqref{eq:rel-th-copt}、LOLP/LOLE/LOLF/LOLD &
  \implfull{src/reliability/reliability\_assessment.cpp:run\_frequency\_duration\_analysis}（独立两状态机组、恒定峰荷；容量不分箱，缺 MTTR 时不伪造 LOLF/LOLD）\\
串并联解析约简~\eqref{eq:rel-th-series}--\eqref{eq:rel-th-parallel} &
  \implfull{src/reliability/reliability\_assessment.cpp:reduce\_series\_frequency\_duration / reduce\_parallel\_frequency\_duration}（独立两状态 CTMC）\\
一般物理网络最小割集解析枚举 &
  \implfull{src/reliability/reliability\_assessment.cpp:evaluate\_physical\_network\_reliability}（成功路径容斥精确可用率、极小击中集、稳定 ID 输出）\\
IEEE 1366 配电指标 &
  \implfull{include/hacdcpf/reliability/reliability\_assessment.hpp:DistributionIndices}（需客户数据）\\
VaR/CVaR 尾部风险 &
  \implfull{include/hacdcpf/reliability/reliability\_assessment.hpp:TailRiskMetrics}（经验 ES 对 VaR 原子作分数权重）\\
混合 AC/DC 有功网络与直流侧削减 &
  \implfull{src/reliability/reliability\_assessment.cpp:evaluate\_hybrid\_fmea\_network\_lp}（\texttt{hybrid-acdc-network-lp}）\\
网络约束状态后果模型~\eqref{eq:rel-rts-edge-set}--\eqref{eq:rel-rts-eens-upper-bound} &
  \implfull{src/optimal\_power\_flow/dc\_opf.cpp:solve\_dc\_opf 与 src/reliability/reliability\_assessment.cpp:evaluate\_state / evaluate\_hybrid\_fmea\_network\_lp}（唯一物理边、逐分量角度参考、$P_g^{\min}=0$、$p^{gc}$、两阶段字典序及规范坐标后验可行性证书）\\
\end{longtable}
\end{center}

\subsection{参考文献}
\begin{thebibliography}{9}\small
\bibitem{rel:billinton-allan} R.~Billinton and R.~N. Allan, \emph{Reliability Evaluation of Power
  Systems}, 2nd ed. New York: Plenum Press, 1996.
\bibitem{rel:billinton-li} R.~Billinton and W.~Li, \emph{Reliability Assessment of Electric Power
  Systems Using Monte Carlo Methods}. New York: Plenum Press, 1994.
\bibitem{rel:rts79} IEEE APM Subcommittee, ``IEEE Reliability Test System,'' \emph{IEEE Trans. Power
  App. Syst.}, vol.~PAS-98, no.~6, pp.~2047--2054, 1979.
\bibitem{rel:rts96} C.~Grigg \emph{et al.}, ``The IEEE Reliability Test System--1996,'' \emph{IEEE
  Trans. Power Syst.}, vol.~14, no.~3, pp.~1010--1020, 1999.
\bibitem{rel:ieee1366} IEEE Std 1366-2012, \emph{IEEE Guide for Electric Power Distribution
  Reliability Indices}. New York: IEEE, 2012.
\bibitem{rel:rockafellar} R.~T. Rockafellar and S.~Uryasev, ``Optimization of conditional
  value-at-risk,'' \emph{Journal of Risk}, vol.~2, no.~3, pp.~21--41, 2000.
\bibitem{rel:zimmerman} R.~D. Zimmerman, C.~E. Murillo-S\'anchez, and R.~J. Thomas,
  ``MATPOWER: Steady-state operations, planning, and analysis tools for power systems
  research and education,'' \emph{IEEE Trans. Power Syst.}, vol.~26, no.~1, pp.~12--19, 2011.
\end{thebibliography}
